% !TEX root = ../main.tex
\lecture{5}{2026-10-12}{Warenkonten}

\section{Warenkonten}

\subsection{Grundproblematik}
Die Lagerung von Waren stellt eine zeitliche Trennung von Produktion (Aufwand) und Realisierung (Ertrag) dar.

Waren zukünftigen Nutzenszufluss aus Verkauf, darum Aktiven
\begin{definition}
    [Einstandspreis] Kosten, die anfallen, bevor und bis die Ware im Unternehmen ist.
    \begin{itemize}
        \item Einkaufspreis
        \item Fracht- und Zollkosten
    \end{itemize}
\end{definition}

Einstandspreise werden in der Bilanz im Warenlagerkonto verbucht. Erst nach dem Verkauf fallen diese als Aufwände an.
Die Kosten der Verweilung der Ware im Unternehmen werden laufend als Aufwände verbucht.


\subsection{Methoden zur Führung der Warenkonten}

\subsubsection{Exakte Methode}
Bei Unternehmen, deren Haupttätigkeit auf Kauf/Verkauf von Waren beruht.
Idee basiert auf die Darstellung der Aktuellen physischen Bewegung der Waren und  Gewinn jederzeit ermittelbar.
Drei Konten werden laufend geführt:
\begin{itemize}
    \item Warenlager (in Einstandspreise)
    \item Warenaufwand (in Einstandspreise)
    \item Warenertrag (in Verkaufspreise) Soll - und Haben +
\end{itemize}

\subsubsection{Praktiker Methode}
Für Organisationen, denen den Warenverkauf keine wesentliche Tätigkeit ist.
Gleichen drei Konten werden verbucht, jedoch nur an Periodenende stimmen die Bezeichnungen.
Es wird vorausgesetzt, dass jede Ware unmittelbar verkauft wird.
Unterjährig werden die Buchungen nicht stimmen, weil durchschnittlich erfasst worden
Prinzip:

\begin{itemize}
    \item \textbf{Wareneinkäufe} als Warenaufwand verbucht
    \item Idee: Kontoführung so einfach wie möglich halten.
\end{itemize}

Korrekturbuchung: am Ende der Periode müssen Korrekturen geführt werden, falls weniger Waren als gekauft verkauft wurden.
\pagebreak
\subsection{Ermittlung des Inventarbestands}

Problematik: Unterjährige Einstandspreise können sich verändern (nach Konjunkur, Lieferanten etc.)
\begin{itemize}
    \item Am zuverlässigsten wäre der exakte Nachweis (aber NICHT IMMER MÖGLICH)
    \item Lösung also:
    \begin{itemize}
        \item Annahme der Zeitlichen Abfolge im Unternehmen in der Periode.
    \end{itemize}
\end{itemize}

\subsubsection{First-In-First-Out (FIFO) }
Nimmt an, dass die zuerst gekauften Waren als Erstes Verkauft werden. \footnote{Wie Fliessband} Die Warenpreise im Lager sind also aktueller, nur die neusten noch übrig.
\subsubsection{Last-In-First-Out (LILO)}
Nimmt an, dass die zuletzt gekauften Waren als Erstes Verkauft werden. \footnote{Platestack}
Von OR erlaubt, aber von internationalen Accounting-Standards verboten, wegen grosser Manipulationsraum aufgrund von alten Warenpreise.
\subsubsection{Durschnittspreis Methode}
\textbf{Nacharbeiten}